\newif\ifglobal

%%%%%%%%%%%%%%%%%%%%%%%
%%% Compile to view %%%
%%%%%%%%%%%%%%%%%%%%%%%

%%%%%%%%%%%%%%%%%%%%%%%%%%%%%%%%%%%%%%%%%%%%%%%%%%%
%%% Readme:                                     %%%
%%% https://www.overleaf.com/read/bwdjvfjksvjy/ %%%
%%%%%%%%%%%%%%%%%%%%%%%%%%%%%%%%%%%%%%%%%%%%%%%%%%%

% >>> M?: marks a module setting number or important notice

% >>> M4: Set {./relative-path} to external files for this module <<<
\def \modulefiles {./21-DBGASSIST}
% To add an external file, use:
% (Replace '---' with actual filename)
% '\input{\modulefiles/tables/---.tex}' for tables
% '\includegraphics{\modulefiles/figures/---}' for figures
% '\subfile{\modulefiles/---}' for register files


% >>> M1: This module is in global project? <<<
% 'YES' - uncomment; 'NO' - comment out
\globaltrue

\ifglobal
    \documentclass[main\_\_EN.tex]{subfiles}
   
\else
    \documentclass[openany, 10pt]{book}

    % >>> M2: In ./00-shared/preamble-trm-module, also find
    %         settings for visibility of todo notes, watermark <<<
    \usepackage{preamble-trm-module}

    % >>> M3: Temporary labels in English,
    %         you can add your own labels there <<<
    \myexternaldocument{./00-shared/temp-labels-en}

\fi %\ifglobal


\setcounter{secnumdepth}{4}


% >>> M5: If this module needs any updates in
% './00-shared/preamble-trm-module.sty'
% (include additional package, etc.), contact your module PM
% or create the file './00-shared/changelog.tex' make the
% necessary updates yourself and document them in this file <<<


\begin{document}


% Sets language into which pre-defined bits of text are translated
\selectlanguage{english}

% Do not delete this block
\ifglobal
% Add the following front matter if
% this module is outside of global project
\else
    \InputIfFileExists{readme}{}{}
    {\let\clearpage\relax \listoftodos}
    {\let\clearpage\relax \listofdonetodos}
    \tableofcontents
    %\listoftables
    %\listoffigures
\fi










%%%%%%%%%%%%%%%%%%%%%%%%%
%%% TRM Module Begins %%%
%%%%%%%%%%%%%%%%%%%%%%%%%

\chapter{Debug Assistant (ASSIST\_DEBUG)}
\label{mod:dbgassist}
\hypertarget{dbgassist}{}


\section{Overview}

Debug Assistant is an auxiliary module that features a set of functions to help locate bugs and issues during software debugging.

\section{Features}
\label{sec:dbgassist-features}

\begin{itemize}
    \item {\bfseries Read/write monitoring}: Monitors whether the High-Performance CPU (HP CPU) bus reads from or writes to a specified memory address space. A detected read or write in the monitored address space will trigger an interrupt.
    \item {\bfseries Stack pointer (SP) monitoring}: Monitors whether the SP exceeds the specified address space. A bounds violation will trigger an interrupt.
    \item {\bfseries Program counter (PC) logging}: Records PC value. The developer can get the last PC value at the most recent HP CPU reset.
    \item{\bfseries Bus access logging}: Records the information about bus access. When the HP CPU, LP CPU, or Direct Memory Access controller (DMA) writes a specified value, the Debug Assistant module will record the data type, address of this write operation, and additionally the PC value when the write is performed by the HP CPU, and push such information to the HP SRAM.
\end{itemize}



\section{Functional Description}
\label{sec:dbgassist-func-descr}

\subsection{Region Read/Write Monitoring}

The Debug Assistant module can monitor reads/writes performed by the HP CPU over Data bus and Peripheral bus in a certain address space, i.e., memory region. Whenever the bus reads or writes in the specified address space, an interrupt will be triggered. The Data bus can monitor two memory regions (assuming they are region 0 and region 1, defined by developers' needs) at the same time, and so can Peripheral Bus.

\subsection{SP Monitoring}

The Debug Assistant module can monitor the SP so as to prevent stack overflow or erroneous push/pop. When the stack pointer exceeds the minimum or maximum threshold, the module will record the PC pointer and generate an interrupt. The threshold is configured by software.

\subsection{PC Logging}

In some cases, software developers want to know the PC at the last HP CPU reset. For instance, when the program is stuck and can only be reset, the developer may want to know where the program got stuck in order to debug. The Debug Assistant module can record the PC at the last HP CPU reset, which can be then read for software debugging. 

\subsection{CPU/DMA Bus Access Logging}

The Debug Assistant module can record the information about the HP CPU Data bus's, LP CPU bus's, and DMA bus's write behaviors in real time. When a write operation occurs in or a specific value is written to a specified address space, the module will record the bus type, the address, PC (only when the write is performed by the HP CPU will PC be recorded), and other information, and then store the data in the HP SRAM in a certain format.

\section{Recommended Operation}
\label{sec:dbgassist-recommended-operation}

\subsection{Region Monitoring and SP Monitoring Configuration}

The Debug Assistant module can monitor reads and writes performed by the HP CPU's Data bus and Peripheral bus. Two memory regions on each bus can be monitored at the same time. All the monitoring modes supported by the Debug Assistant module are listed below:

\begin{itemize}
    \item Monitoring of the read/write operations performed by Data bus
        \begin{itemize}
            \item Data bus reads in region 0
            \item Data bus writes in region 0
            \item Data bus reads in region 1
            \item Data bus writes in region 1
        \end{itemize}    
    \item Monitoring of the read/write operations performed by Peripheral bus
        \begin{itemize}
            \item Peripheral bus reads in region 0
            \item Peripheral bus writes in region 0
            \item Peripheral bus reads in region 1
            \item Peripheral bus writes in region 1
        \end{itemize}
    \item  Monitoring of exceeding the SP bounds
        \begin{itemize}
            \item SP exceeds the upper bound address
            \item SP exceeds the lower bound address
        \end{itemize}
\end{itemize}



The configuration process for region monitoring and SP monitoring is as follows:

\begin{enumerate}
    \item Configure monitored region and SP threshold.
        \begin{itemize}
            \item Configure Data bus region 0 with \hyperref[regdesc:ASSISTDEBUGCORE0AREADRAM00MINREG]{ASSIST\_DEBUG\_CORE\_0\_AREA\_DRAM0\_0\_MIN\_REG} and\\ \hyperref[regdesc:ASSISTDEBUGCORE0AREADRAM00MAXREG]{ASSIST\_DEBUG\_CORE\_0\_AREA\_DRAM0\_0\_MAX\_REG}. 
            \item Configure Data bus region 1 with \hyperref[regdesc:ASSISTDEBUGCORE0AREADRAM01MINREG]{ASSIST\_DEBUG\_CORE\_0\_AREA\_DRAM0\_1\_MIN\_REG} and\\ \hyperref[regdesc:ASSISTDEBUGCORE0AREADRAM01MAXREG]{ASSIST\_DEBUG\_CORE\_0\_AREA\_DRAM0\_1\_MAX\_REG}.
            \item Configure Peripheral bus region 0 with \hyperref[regdesc:ASSISTDEBUGCORE0AREAPIF0MINREG]{ASSIST\_DEBUG\_CORE\_0\_AREA\_PIF\_0\_MIN\_REG} and\\ \hyperref[regdesc:ASSISTDEBUGCORE0AREAPIF0MAXREG]{ASSIST\_DEBUG\_CORE\_0\_AREA\_PIF\_0\_MAX\_REG}.
            \item Configure Peripheral bus region 1 with \hyperref[regdesc:ASSISTDEBUGCORE0AREAPIF1MINREG]{ASSIST\_DEBUG\_CORE\_0\_AREA\_PIF\_1\_MIN\_REG} and\\ \hyperref[regdesc:ASSISTDEBUGCORE0AREAPIF1MAXREG]{ASSIST\_DEBUG\_CORE\_0\_AREA\_PIF\_1\_MAX\_REG}.
            \item Configure SP threshold with \hyperref[regdesc:ASSISTDEBUGCORE0SPMINREG]{ASSIST\_DEBUG\_CORE\_0\_SP\_MIN\_REG} and\\ \hyperref[regdesc:ASSISTDEBUGCORE0SPMAXREG]{ASSIST\_DEBUG\_CORE\_0\_SP\_MAX\_REG}.
        \end{itemize}
    
    \item Configure interrupts.
    \begin{itemize}
        \item Configure \hyperref[regdesc:ASSISTDEBUGCORE0INTRENAREG]{ASSIST\_DEBUG\_CORE\_0\_INTR\_ENA\_REG} to enable the interrupt of a monitoring mode. 
        \item Configure \hyperref[regdesc:ASSISTDEBUGCORE0INTRRAWREG]{ASSIST\_DEBUG\_CORE\_0\_INTR\_RAW\_REG} to get the interrupt status of a monitoring mode. 
        \item Configure \hyperref[regdesc:ASSISTDEBUGCORE0INTRCLRREG]{ASSIST\_DEBUG\_CORE\_0\_INTR\_CLR\_REG} to clear the interrupt of a monitoring mode.
    \end{itemize}
    \item Configure \hyperref[regdesc:ASSISTDEBUGCORE0MONTRENAREG]{ASSIST\_DEBUG\_CORE\_0\_MONTR\_ENA\_REG} to enable the monitoring mode(s). Various monitoring modes can be enabled at the same time.
\end{enumerate}

Assuming that Debug Assistant module needs to monitor whether Data bus has written to [A \verb+~+ B] address space, the user can enable monitoring in either Data bus region 0 or region 1. The following configuration process is based on region 0:

\begin{enumerate}
    \item Configure \hyperref[fielddesc:ASSISTDEBUGCORE0RCDPDEBUGEN]{ASSIST\_DEBUG\_CORE\_0\_RCD\_PDEBUGEN} to 1 to enable HP CPU to update the PC signals to the Debug Assistant module.
    \item Configure \hyperref[regdesc:ASSISTDEBUGCORE0AREADRAM00MINREG]{ASSIST\_DEBUG\_CORE\_0\_AREA\_DRAM0\_0\_MIN\_REG} to Address A.
    \item Configure \hyperref[regdesc:ASSISTDEBUGCORE0AREADRAM00MAXREG]{ASSIST\_DEBUG\_CORE\_0\_AREA\_DRAM0\_0\_MAX\_REG} to Address B.
    \item Configure \hyperref[regdesc:ASSISTDEBUGCORE0INTRENAREG]{ASSIST\_DEBUG\_CORE\_0\_INTR\_ENA\_REG} bit[1] to enable the interrupt for write operations by Data bus in region 0.
    \item Configure \hyperref[regdesc:ASSISTDEBUGCORE0MONTRENAREG]{ASSIST\_DEBUG\_CORE\_0\_MONTR\_ENA\_REG} bit[1] to enable monitoring write operations by Data bus in region 0.
    \item Configure interrupt matrix to map ASSIST\_DEBUG\_INT into HP CPU interrupt (please refer to Chapter \ref{mod:intmtrx} \textit{\nameref{mod:intmtrx}}).
    \item After the interrupt is triggered:
    
    \begin{itemize}
        \item Read \hyperref[regdesc:ASSISTDEBUGCORE0INTRRAWREG]{ASSIST\_DEBUG\_CORE\_0\_INTR\_RAW\_REG} to learn which operation triggered the interrupt.
        
        \item If the interrupt is triggered by region monitoring, read \hyperref[regdesc:ASSISTDEBUGCORE0AREAPCREG]{ASSIST\_DEBUG\_CORE\_0\_AREA\_PC\_REG} for the PC value, and \hyperref[regdesc:ASSISTDEBUGCORE0AREASPREG]{ASSIST\_DEBUG\_CORE\_0\_AREA\_SP\_REG} for the SP.
        
        \item If the interrupt is triggered by stack monitoring, read \hyperref[regdesc:ASSISTDEBUGCORE0SPPCREG]{ASSIST\_DEBUG\_CORE\_0\_SP\_PC\_REG} for the PC value.
        
        \item Write 1 to the corresponding bits of \hyperref[regdesc:ASSISTDEBUGCORE0INTRRAWREG]{ASSIST\_DEBUG\_CORE\_0\_INTR\_RAW\_REG} to clear the interrupts.
    \end{itemize}

\end{enumerate}

\subsection{PC Logging Configuration}


Configure \hyperref[fielddesc:ASSISTDEBUGCORE0RCDPDEBUGEN]{ASSIST\_DEBUG\_CORE\_0\_RCD\_PDEBUGEN} to 1 to enable HP CPU to update the PC signals to the Debug Assistant module. If \hyperref[fielddesc:ASSISTDEBUGCORE0RCDRECORDEN]{ASSIST\_DEBUG\_CORE\_0\_RCD\_RECORDEN} is also configured to 1, \\\hyperref[regdesc:ASSISTDEBUGCORE0RCDPDEBUGPCREG]{ASSIST\_DEBUG\_CORE\_0\_RCD\_PDEBUGPC\_REG} will record the HP CPU's PC signal and \hyperref[regdesc:ASSISTDEBUGCORE0RCDPDEBUGSPREG]{ASSIST\_DEBUG\_CORE\_0\_RCD\_PDEBUGSP\_REG} will record the SP value. Otherwise, the two registers keep the original values.

When the CPU resets, \hyperref[regdesc:ASSISTDEBUGCORE0RCDENREG]{ASSIST\_DEBUG\_CORE\_0\_RCD\_EN\_REG} will reset, while \\ \hyperref[regdesc:ASSISTDEBUGCORE0RCDPDEBUGPCREG]{ASSIST\_DEBUG\_CORE\_0\_RCD\_PDEBUGPC\_REG}  and \hyperref[regdesc:ASSISTDEBUGCORE0RCDPDEBUGSPREG]{ASSIST\_DEBUG\_CORE\_0\_RCD\_PDEBUGSP\_REG} will not. Therefore, the two registers will keep the PC value and SP value at the CPU reset.

\subsection{CPU/DMA Bus Access Logging Configuration}

The configuration process for CPU/DMA bus access logging is described below.

\begin{enumerate}

\item Configure monitored address space.
      \begin{itemize}
         \item Configure \hyperref[regdesc:MEMMONITORLOGMINREG]{MEM\_MONITOR\_LOG\_MIN\_REG} and \hyperref[regdesc:MEMMONITORLOGMAXREG]{MEM\_MONITOR\_LOG\_MAX\_REG} to specify monitored address space.
      \end{itemize}

\item Configure the monitoring mode with \hyperref[fielddesc:MEMMONITORLOGMODE]{MEM\_MONITOR\_LOG\_MODE}:
       \begin{itemize}
           \item write monitoring (whether the bus has write operations)
           \item word monitoring (whether the bus writes a specific word)
           \item halfword monitoring (whether the bus writes a specific halfword)
           \item byte monitoring (whether the bus writes a specific byte)
       \end{itemize}
       
\item Configure the specific values to be monitored.
       \begin{itemize}
            \item In word monitoring mode, \hyperref[regdesc:MEMMONITORLOGCHECKDATAREG]{MEM\_MONITOR\_LOG\_CHECK\_DATA\_REG} specifies the monitored word.
            \item In halfword monitoring mode, \hyperref[regdesc:MEMMONITORLOGCHECKDATAREG]{MEM\_MONITOR\_LOG\_CHECK\_DATA\_REG}[15:0] specifies the monitored halfword.
            \item In byte monitoring mode, \hyperref[regdesc:MEMMONITORLOGCHECKDATAREG]{MEM\_MONITOR\_LOG\_CHECK\_DATA\_REG}[7:0] specifies the monitored byte.
            \item \hyperref[regdesc:MEMMONITORLOGDATAMASKREG]{MEM\_MONITOR\_LOG\_DATA\_MASK\_REG} is used to mask the byte specified in \hyperref[regdesc:MEMMONITORLOGCHECKDATAREG]{MEM\_MONITOR\_LOG\_CHECK\_DATA\_REG}. A masked byte can be any value. For example, in word monitoring, if \hyperref[regdesc:MEMMONITORLOGCHECKDATAREG]{MEM\_MONITOR\_LOG\_CHECK\_DATA\_REG} is configured to 0x01020304 and \hyperref[regdesc:MEMMONITORLOGDATAMASKREG]{MEM\_MONITOR\_LOG\_DATA\_MASK\_REG} is configured to 0x1, then any writes of the data matching the 0x010203XX pattern by the bus will be recorded.
       \end{itemize}
       
\item Configure the storage space for recorded data.
        \begin{itemize}
            \item \hyperref[regdesc:MEMMONITORLOGMEMSTARTREG]{MEM\_MONITOR\_LOG\_MEM\_START\_REG} and \hyperref[regdesc:MEMMONITORLOGMEMENDREG]{MEM\_MONITOR\_LOG\_MEM\_END\_REG} specify the storage space for recorded data. The storage space must be in the range of 0x4080\_0000 \verb+~+ 0x4087\_FFFF.
            \item Set \hyperref[regdesc:MEMMONITORLOGMEMADDRUPDATEREG]{MEM\_MONITOR\_LOG\_MEM\_ADDR\_UPDATE\_REG} to update the value in \hyperref[regdesc:MEMMONITORLOGMEMSTARTREG]{MEM\_MONITOR\_LOG\_MEM\_START\_REG} to \hyperref[regdesc:MEMMONITORLOGMEMCURRENTADDRREG]{MEM\_MONITOR\_LOG\_MEM\_\\CURRENT\_ADDR\_REG}.
            \item Configure the permission for the Debug Assistant module to access the internal HP SRAM. Only when the access permission is enabled can the Debug Assistant module access the internal HP SRAM. For more information, please refer to Chapter \ref{mod:permctrl} \textit{\nameref{mod:permctrl}}.
        \end{itemize}

\item Configure the writing mode for the recorded data: loop mode or non-loop mode.
        \begin{itemize}
            \item In loop mode, writing to the specified address space is performed in loops. When writing reaches the end address, it will return to the starting address and continue, overwriting the previously recorded data. Set \hyperref[fielddesc:MEMMONITORLOGMEMLOOPENABLE]{MEM\_MONITOR\_LOG\_MEM\_LOOP\_ENABLE} to enable loop mode.\\
           For example, there are 10 write operations (1 \verb+~+ 10) to address space 0 \verb+~+ 4 during bus access. After the 5th operation writes to address 4, the 6th operation will start writing from address 0. The 6th to 10th operations will overwrite the previous data written by the 1th to 5th operations.
            \item In non-loop mode, when writing reaches the end address, it will stop at the end address and dump the remaining data, not overwriting the previously recorded data. Clear \hyperref[fielddesc:MEMMONITORLOGMEMLOOPENABLE]{MEM\_MONITOR\_LOG\_MEM\_LOOP\_ENABLE} to use non-loop mode.\\
            For example, there are 10 write operations (1 \verb+~+ 10) to address space 0 \verb+~+ 4 during bus access. After the 5th operation writes to address 4, the 6th to 10th write operations will stop at address 4 and will not be performed any more. Therefore, the address 0 \verb+~+ 4 stores the values written by the 1 \verb+~+ 5 operations and the values of the 6 \verb+~+ 10 operations are dumped.
        \end{itemize}

\item Configure bus enable registers.
        \begin{itemize}
            \item Enable HP CPU, LP CPU, or DMA bus access logging with \hyperref[fielddesc:MEMMONITORLOGENA]{MEM\_MONITOR\_LOG\_ENA}. They can be enabled at the same time.
        \end{itemize}
\end{enumerate}

The Debug Assistant module first writes the recorded data to an internal buffer, and then fetches the data from the buffer and writes it to the configured memory space. When the monitored behaviors are triggered continuously, the generated recording packets may fully occupy the buffer, making it unable to take any incoming packets. At this time, the module dumps these incoming packets and buffers a LOST packet instead before the buffer reaches its capacity. However, the bus type and the number of these dumped packets are unknown.

When bus access logging is finished, the recorded data can be read from memory for decoding. The recorded data is in four packet formats, namely HP CPU packet (corresponding to HP CPU Data bus), LP CPU packet (corresponding to LP CPU bus), DMA packet (corresponding to DMA bus), and LOST packet. The packet formats are shown in Table \ref{tab:hp_cpu_package}, \ref{tab:lp_cpu_package}, \ref{tab:dma_package}, and \ref{tab:lost_package}.

    \begin{longtable}[c]{|L{3cm}|L{3cm}|L{3cm}|L{3cm}|l|}
    \caption{HP CPU Packet Format}\label{tab:hp_cpu_package}
    \\\hline
    \rowcolor{lightgray}
    {\bfseries Bit[63:34]} & {\bfseries Bit[33:32]} & {\bfseries Bit[31:4]} & {\bfseries Bit[3:2]} & {\bfseries Bit[1:0]} \\\hline
    pc\_offset & anchored(2) & addr\_offset & format & anchored(1) \\\hline
    \end{longtable}
    
    \begin{longtable}[c]{|L{3cm}|L{3cm}|l|}
    \caption{LP CPU Packet Format}\label{tab:lp_cpu_package}
    \\\hline
    \rowcolor{lightgray}
    {\bfseries Bit[31:4]} & {\bfseries Bit[3:2]} & {\bfseries Bit[1:0]} \\\hline
    addr\_offset & format & anchored(1) \\\hline
    \end{longtable}
    
    \begin{longtable}[c]{|L{3cm}|L{2cm}|L{2cm}|l|}
    \caption{DMA Packet Format}\label{tab:dma_package}
    \\\hline
    \rowcolor{lightgray}
    {\bfseries Bit[31:9]} & {\bfseries Bit[8:4]} & {\bfseries Bit[3:2]} & {\bfseries Bit[1:0]} \\\hline
     addr\_offset & dma\_source & format & anchored(1)  \\\hline
    \end{longtable}
    
    \begin{longtable}[c]{|L{3cm}|L{3cm}|l|}
    \caption{LOST Packet Format}\label{tab:lost_package}
    \\\hline
    \rowcolor{lightgray}
    {\bfseries Bit[31:4]} & {\bfseries Bit[3:2]} & {\bfseries Bit[1:0]} \\\hline
    reserved & format & anchored(1) \\\hline
    \end{longtable}

It can be seen from the data packet formats that the HP CPU packet size is 64 bits, LP CPU packet 32 bits, DMA packet size 32 bits, and LOST packet 32 bits. These packets contain the following fields:

\begin{itemize}
    \item {\bfseries format} – the packet type. 0: HP CPU packet; 1: DMA packet; 2: LP CPU packet; 3: LOST packet.
    \item {\bfseries pc\_offset} - the offset of the PC register at the time of access. Actual PC = pc\_offset + 0x4000\_0000.
    \item {\bfseries addr\_offset} - the address offset of a write operation. Actual address = addr\_offset + \hyperref[regdesc:MEMMONITORLOGMINREG]{MEM\_MONITOR\_LOG\_MIN\_REG}.
    \item {\bfseries dma\_source} - the source of DMA access. Refer to Table \ref{tab:dma_source}. For more information on the values 16 \verb+~+ 31 in the table, please refer to \ref{mod:dma} \textit{\nameref{mod:dma}}.
    \item {\bfseries anchored} - the location of the 32 bits in the data packet. 1 indicates the lower 32 bits. 2 indicates the higher 32 bits.
    
        \begin{longtable}[c]{|l|m{8cm}|}
        \caption{DMA Access Source}\label{tab:dma_source}
        \\\hline
        \rowcolor{lightgray}
        {\bfseries Value} & {\bfseries Source} \\\hline
        \endfirsthead \hline \rowcolor{lightgray}
        {\bfseries Value} & {\bfseries Source} \\\hline
        \endhead \endfoot \hline \endlastfoot
        0 & HP CPU \\\hline
        1 & LP CPU \\\hline
        2 & reserved \\\hline
        3 & SDIO\_SLV \\\hline
        4 & reserved \\\hline
        5 & MEM\_MONITOR \\\hline
        6 & TRACE \\\hline 
        7 \verb+~+ 15 & reserved \\\hline
        16 \verb+~+ 31 & See the peripherals corresponding to the values 0 \verb+~+ 15 in Chapter \ref{mod:dma} \textit{\nameref{mod:dma}} > Table \ref{tab:gdma-peri-sel} \textit{\nameref{tab:gdma-peri-sel}}. For example, the source corresponding to the value 16 is the peripheral corresponding to the value 0 in that table, the source corresponding to 17 is the peripheral corresponding to 1 in that table, and etc.\\\hline
        \end{longtable}
\end{itemize}    


The internal buffer of the module is 32 bits wide. When the HP CPU, LP CPU, or DMA bus access logging are all enabled at the same time and the record data is generated at the same time, the DMA data packets are first buffered, then the HP CPU packets, and finally the LP CPU packets. This priority of buffering packets also applies to the case where only two types of packets are generated at the same time. The Debug Assistant will automatically fetch the buffered data and store it in 32-bit data width into the specified memory space.

In loop mode, data looping several times in the storage memory may cause residual data, which can interfere with packet parsing. For example, the lower 32 bits of a HP CPU packet are overwritten, thus making its higher 32 bits residual data. Therefore, users need to filter out the possible residual data in order to determine the starting position of the first valid packet with \hyperref[regdesc:MEMMONITORLOGMEMCURRENTADDRREG]{MEM\_MONITOR\_LOG\_MEM\_CURRENT\_ADDR\_REG}. Once the starting position of the packet is identified, check the anchored bit value of the packet. If it is 1, the data will be retained. If it is 2, it will be dumped. 

The process of packet parsing is described below:

\begin{itemize}
    \item Determine whether there is a data overflow with \hyperref[fielddesc:MEMMONITORLOGMEMFULLFLAG]{MEM\_MONITOR\_LOG\_MEM\_FULL\_FLAG}. If no, the address space to read is \hyperref[regdesc:MEMMONITORLOGMEMSTARTREG]{MEM\_MONITOR\_LOG\_MEM\_START\_REG}  \verb+~+ \hyperref[regdesc:MEMMONITORLOGMEMCURRENTADDRREG]{MEM\_MONITOR\_LOG\_MEM\_\\CURRENT\_ADDR\_REG} - 4. If yes and the loop mode is enabled, the address space is \hyperref[regdesc:MEMMONITORLOGMEMCURRENTADDRREG]{MEM\_MONITOR\_\\LOG\_MEM\_CURRENT\_ADDR\_REG} \verb+~+ \hyperref[regdesc:MEMMONITORLOGMEMENDREG]{MEM\_MONITOR\_LOG\_MEM\_END\_REG} and
    \hyperref[regdesc:MEMMONITORLOGMEMSTARTREG]{MEM\_MONITOR\_\\LOG\_MEM\_START\_REG}  \verb+~+ \hyperref[regdesc:MEMMONITORLOGMEMCURRENTADDRREG]{MEM\_MONITOR\_LOG\_MEM\_CURRENT\_ADDR\_REG} - 4. If yes and loop mode is not enabled, the address space is \hyperref[regdesc:MEMMONITORLOGMEMSTARTREG]{MEM\_MONITOR\_LOG\_MEM\_START\_REG} \verb+~+ \hyperref[regdesc:MEMMONITORLOGMEMENDREG]{MEM\_MONITOR\_\\LOG\_MEM\_END\_REG}.
    
    \item Read and parse data from the starting address. Read 32 bits each time.
\end{itemize}

After packet parsing is completed, clear the \hyperref[fielddesc:MEMMONITORLOGMEMFULLFLAG]{MEM\_MONITOR\_LOG\_MEM\_FULL\_FLAG} flag bit by setting \hyperref[fielddesc:MEMMONITORCLRLOGMEMFULLFLAG]{MEM\_MONITOR\_CLR\_LOG\_MEM\_FULL\_FLAG}.

%%%%%%%%%%%%%%%%%%%%%%%%
%%% Register Summary %%%
%%%%%%%%%%%%%%%%%%%%%%%%

\clearpage
\hypertarget{dbgassist-reg-summ}{}
\section{Register Summary}
\label{sec:dbgassist-reg-summ}

The addresses of \textbf{bus logging configuration registers} (see \ref{sec:dbgassist-mem-monitor-reg-summ}) in this section are relative to the \textbf{MEM\_MONITOR} base address. The addresses of other registers (see \ref{sec:dbgassist-bus-monitor-reg-summ}) are relative to the \textbf{ASSIST\_DEBUG} base address. Both base addresses are provided in Table \ref{tab:sysmem-base-address} in Chapter \ref{mod:sysmem} \textit{\nameref{mod:sysmem}}.

The abbreviations given in Column \textbf{Access} are explained in Section \hyperref[glossary-access-types]{\textit{Access Types for Registers}}.

\subsection{Summary of Bus Logging Configuration Registers}
\label{sec:dbgassist-mem-monitor-reg-summ}

\subfile{\modulefiles/21-DBGASSIST-mem-monitor-reg-summ__EN}

\subsection{Summary of Other Registers}
\label{sec:dbgassist-bus-monitor-reg-summ}

\subfile{\modulefiles/21-DBGASSIST-reg-summ__EN}

%%%%%%%%%%%%%%%%%
%%% Registers %%%
%%%%%%%%%%%%%%%%%

% >>> Update the sample text below and insert
%      the info on registers into the subfile <<<

\clearpage
\section{Registers}

The addresses of \textbf{bus logging configuration registers} (see \ref{sec:dbgassist-mem-monitor-reg}) in this section are relative to \textbf{MEM\_MONITOR} base address. The addresses of other registers (see \ref{sec:dbgassist-bus-monitor-reg}) are relative to the \textbf{ASSIST\_DEBUG} base address. Both base addresses are provided in Table \ref{tab:sysmem-base-address} in Chapter \ref{mod:sysmem} \textit{\nameref{mod:sysmem}}.

\subsection{Bus Logging Configuration Registers}
\label{sec:dbgassist-mem-monitor-reg}

\subfile{\modulefiles/21-DBGASSIST-mem-monitor-reg__EN}

\subsection{Other Registers}
\label{sec:dbgassist-bus-monitor-reg}

\subfile{\modulefiles/21-DBGASSIST-reg__EN}


\end{document}
